%! AP Statistics — Unit 4, Lesson 4.1 — Student Packet Cover Sheet
\documentclass[10pt]{article}
\usepackage{apstats-article}
\usepackage{apstats-boxes}
\usepackage{ltablex}
\keepXColumns

\begin{document}

% Full-bleed navy banner
\begin{tikzpicture}[remember picture, overlay]
  \fill[navy] (current page.north west)
    rectangle ([yshift=-0.9in]current page.north east);
\end{tikzpicture}
\vspace{-0.8in}

\begin{center}
  {\color{white}\textbf{\LARGE AP Statistics}}\\[4pt]
  {\color{skymid}\large\textbf{Unit 4: Probability, Random Variables, and Probability Distributions}}\\[6pt]
  {\color{white}\textbf{\large Lesson 4.1 \quad Introducing Statistics: Random and Non-Random Patterns?}}
\end{center}

\vspace{0.22in}
\namedateperiod
\vspace{0.18in}

\begin{learningtargetbox}
\small
\begin{enumerate}[leftmargin=1.5em, itemsep=5pt]
  \item \ldots{} identify patterns in data and write statistical questions those patterns suggest.
  \item \ldots{} explain why a pattern alone does not confirm a non-random process.
  \item \ldots{} recognize that random processes can produce streaks, clusters, and other surprising patterns.
\end{enumerate}
\end{learningtargetbox}

\vspace{0.14in}

\begin{tocbox}
\small
\renewcommand{\arraystretch}{1.5}
\begin{tabularx}{\linewidth}{@{} c l X r @{}}
\rowcolor{sky}
\textbf{\#} & \textbf{Component} & \textbf{Description} & \textbf{Score} \\
\midrule
1 & Warm-Up       & Spiral review + pattern recognition thinking       & \blank{1.2cm} \\
2 & Guided Notes  & Random vs.\ non-random patterns; statistical questions & \blank{1.2cm} \\
3 & Group Activity & Evaluating real-world patterns for randomness      & \blank{1.2cm} \\
4 & Exit Ticket   & Die-rolling sequence analysis                      & \blank{1.2cm} \\
5 & Homework      & Four pattern scenarios; reflection                  & \blank{1.2cm} \\
\midrule
  & \multicolumn{2}{r}{\textbf{Total}} & \blank{1.2cm} \\
\end{tabularx}
\end{tocbox}

\vspace{0.14in}

\begin{remindbox}
\small
The \textbf{most important idea in Lesson 4.1}: a pattern in data is not evidence of a
non-random cause. Random processes naturally produce streaks and clusters.

\vspace{6pt}
\begin{tabularx}{\linewidth}{@{} >{\centering\arraybackslash\columncolor{goldbg}}p{0.06\linewidth}
                                    X @{}}
\textbf{\large!} &
\textbf{Patterns are not proof.}\enspace
Even a fair coin can produce 6 heads in a row. Seeing a pattern should inspire a
question, not a conclusion. \\[6pt]
\textbf{\large!} &
\textbf{Our intuition about randomness is often wrong.}\enspace
People expect random sequences to look more ``mixed up'' than they really are; true
random sequences contain more long runs than most people expect. \\[6pt]
\textbf{\large!} &
\textbf{Ask the right question.}\enspace
``Could chance alone produce this?'' is the key statistical question. Answering it
requires probability --- which is what the rest of Unit 4 develops. \\
\end{tabularx}
\end{remindbox}

\end{document}
